\input{_preamble}
\begin{document}
\ex. Alpha.

\ex.
\a. Beta.
\b. Gamma.
\a. Delta.
\b. Epsilon.

\ex.[(x)] Custom.

\ex. Zeta.

% The same optional argument on the glossed shorthand.  It used to be
% impossible: \exg. puts \lx@glosshead in front of the body, so the
% bracket was no longer the first thing the label peek saw, and the label
% came out as the first word of the object line while the example took a
% number of its own.
\exg.[(xx)] Der Hund bellte.\\
     the dog barked\\

% ... and a bracket a SPACE separates from \exg. is NOT a label: it is
% the first word of the object line, which is what "[" is doing among the
% openers of \GlossPhantomChars.  This is where \exg. parts company with
% \ex., whose optional argument follows the LaTeX convention of skipping
% spaces because it has no object line to compete with.
\exg. [DP der Hund] bellte.\\
      [DP the  dog] barked\\
\end{document}
